\documentclass[10pt,a4paper]{amsart}
\usepackage[margin=19mm]{geometry}
\usepackage{amsmath,amssymb,mathtools}
\usepackage[scr=rsfs,cal=euler]{mathalfa}
\usepackage{xfrac}
\usepackage[unicode,hidelinks,bookmarks=false]{hyperref}
\newcommand{\ie}{i.e.\ }
\newcommand{\cf}{cf.\ }
\newcommand{\resp}{respectively\ }
\newcommand{\Subsection}{\subsection}
\providecommand{\nobreakdash}{\nobreak\hbox{-}}
\setlength{\emergencystretch}{2em}
\multlinegap0pt
\usepackage{amssymb}
\usepackage[all]{xy}
\usepackage{mathtools}
\usepackage{xcolor}
\newtheorem{theorem}{Theorem}
\newtheorem{prop}{Proposition}
\providecommand{\abs}[1]{\lvert#1\rvert}
\newcommand{\bA}{\mathbb{A}}
\newcommand{\cF}{\mathcal{F}}
\newcommand{\cS}{\mathcal{S}}
\newcommand{\fX}{\mathfrak{X}}
\newcommand{\fa}{\mathfrak{a}}
\newcommand{\rd}{\mathrm{d}}
\title{The global Gan--Gross--Prasad conjecture for Fourier--Jacobi periods on unitary groups I: Coarse expansions of the relative trace formulae}
\author{Paul Boisseau, Weixiao Lu, Hang Xue}
\date{Source: arXiv:2404.07342v3 (CC BY 4.0)}
\begin{document}
\maketitle

\noindent\emph{Source and licence.} The global Gan--Gross--Prasad conjecture for Fourier--Jacobi periods on unitary groups I: Coarse expansions of the relative trace formulae. Version arXiv:2404.07342v3 (CC BY 4.0), distributed by arXiv under the Creative Commons Attribution 4.0 International licence, http://creativecommons.org/licenses/by/4.0/, as recorded in the arXiv licence metadata for this exact version and shown on the abstract page https://arxiv.org/abs/2404.07342v3. Full licence notice: \texttt{licenses/arxiv-2404.07342-CC-BY-4.0.txt}.\par\smallskip

\noindent\textit{Editorial scope.} Counted unit: the article's prop prop:invariant\_I\_chi (source0.tex lines 4362--4371). The article's complete original proof of it is delivered with it (source0.tex lines 4373--4513, one \texttt{proof} environment of 141 lines). No mathematics is written, repaired or re-derived: the statement and the proof are the article's own lines, in the article's own notation, and no retained line is substituted or re-flowed.  Retained uncounted as the source-local context the proof needs: lines 4049--4059; lines 4132--4136.  Nothing was omitted from any retained span, so the package carries no omission record.  No retained line contains a citation, so the wrapper carries no bibliography and no bibliography entry.  No retained line contains a figure environment, an image-inclusion command or drawing code, so no image is delivered and the package declares no asset dependency.  Editorial formatting only, in addition: an AMS \texttt{amsart} preamble that restates the article's own declarations and macros used by the retained lines, namely the environments \texttt{theorem}, \texttt{prop} on separate counters, together with the standard packages those lines need.  The retained environments print in this document as Theorem 1, Proposition 1 in order of appearance, whereas the article prints the same items with the numbering of its own full document.  The complete native source of the article as distributed in the arXiv e-print for this exact version is bundled unchanged as \texttt{sources/157531/source0.tex}.
\begin{theorem} \label{thm:coarse_spectral_expansion_GL}
For $T$ sufficiently positive, the functions $I^T(f_+)$ and $I_{\chi}^T(f_+)$ are the
restrictions of exponential-polynomial functions whose purely polynomial parts
are constants. We denote them by $I(f_+)$ and $I_{\chi}(f_+)$ respectively.
Then $I$ and $I_\chi$ are continuous as distributions on $\cS(G_+(\bA))$
and for any $f_+ \in \cS(G_+(\bA))$ we have
    \[
    \sum_{\chi \in \fX(G)} I_{\chi}(f_+) = I(f_+),
    \]
where the sum is absolutely convergent.
\end{theorem}

    \begin{equation}    \label{eq:Gamma_Q'}
    \widehat{\tau}_{P_{n+1}}(H-X) = \sum_{ \substack{ Q \in \cF \\ Q \supset P }}
    \epsilon_{Q} \widehat{\tau}_{P_{n+1}}^{Q_{n+1}}(H)
    \Gamma'_{Q_{n+1}}(H, X), \quad H, X \in \fa_{P_{n+1}}^{G_{n+1}}.
    \end{equation}

\begin{prop}    \label{prop:invariant_I_chi}
The distribution $I_\chi$ is left $H(\bA)$-invariant and right
$(G'(\bA),\eta_{G'})$-equivariant. More precisely for all $f_+ \in
\cS(G_+(\bA))$, $h \in H(\bA)$ and $g' \in G'(\bA)$ we have
    \[
    I_\chi((h, g') \cdot f_+)= \eta_{n+1}(g') I_\chi(f_+)
    \]
where $((h, g') \cdot f_+)(g,u)=\abs{\det h}_E^{1/2} \mu(\det h)^{-1}
f_+(h^{-1} g g',uh)$.
\end{prop}

\begin{proof}
Fix an $(h_0,g_0') \in H(\bA) \times G'(\bA)$. Let $T \in \fa_{n+1}$ be
sufficiently positive. By~\eqref{eq:Gamma_Q'} again we have
    \begin{align*}
     \widehat{\tau}_{P_{n+1}}(H_{P_{n+1}}(\delta_1 g_1'g_{0,1}^{\prime})
     -T_{P_{n+1}})
     = & \sum_{P \subset Q}  \epsilon_Q \widehat{\tau}_{P_{n+1}}^{Q_{n+1}}
     ( H_{P_{n+1}}(\delta_1 g_1') - T_{P_{n+1}} )\\
     &\Gamma'_{Q_{n+1}} ( H_{Q_{n+1}}(\delta_1 g_1')-T_{Q_{n+1}},H_{Q_{n+1}}
     (\delta_1 g_1') - H_{Q_{n+1}}(\delta_1 g_1' g_{0,1}^{\prime})).
    \end{align*}
Therefore we see that $K_{f_+,\chi}^T(h h_0, g' g_0')$ equals
    \[
    \begin{aligned}
    \sum_{P \in \cF}
    \sum_{P \subset Q} \sum_{ \substack{\gamma \in P_H(F)
    \backslash H(F) \\ \delta \in P_{G'}(F)\backslash G'(F)}}
    & \epsilon_Q \widehat{\tau}_{P_{n+1}}^{Q_{n+1}}
     ( H_{P_{n+1}}(\delta_1 g_1') - T_{P_{n+1}} )
     K_{f_+, P, \chi}(\gamma h h_0,\delta g' g_0')\\
    &\Gamma'_{Q_{n+1}}(H_{Q_{n+1}}(\delta_1 g_1')-T_{Q_{n+1}},H_{Q_{n+1}}
    (\delta_1 g_1')-H_{Q_{n+1}}(\delta_1 g_1' g_{0,1}^{\prime})).
    \end{aligned}
    \]
Note that
    \[
    K_{(h_0, g_0') \cdot f_+, P, \chi}(h,g')=
    K_{f_+, P, \chi}(h h_0, g' g_0').
    \]
Thus by summing over $P \in \cF$ first in  the above expression, we see that
it simplifies to
    \begin{equation}    \label{eq:invariance_GL_1}
    \begin{aligned}
    \sum_{Q \in \cF} \sum_{ \substack{\gamma \in Q_H(F)
    \backslash H(F) \\ \delta \in Q_{G'}(F)\backslash G'(F)}}
    &\Gamma'_{Q_{n+1}}(H_{Q_{n+1}}(\delta_1 g_1')-T_{Q_{n+1}},H_{Q_{n+1}}
    (\delta_1 g_1')-H_{Q_{n+1}}(\delta_1 g_1' g_{0,1}^{\prime}))\\
    &K_{(h_0, g_0') \cdot f_+, \chi}^{Q, T}(\gamma h,\delta g').
    \end{aligned}
    \end{equation}

We integrate $K_{f_+,\chi}^T(h h_0, g' g_0') \eta_{n+1}(g')$ over $[H] \times [G']$.
On the one hand by definition it equals
    \[
    \int_{[H] \times [G']}
    K_{f_+,\chi}^T(hh_0,g'g_0^{\prime}) \eta_{n+1}(g') \rd g' \rd h
    =  \eta_{n+1}(g'_0)  I^T_\chi(f_+).
    \]
On the other hand, by~\eqref{eq:invariance_GL_1},
it equals the sum over all $Q \in \cF$ of the terms
    \begin{equation}    \label{eq:invariant_Q_GL}
    \begin{aligned}
    \int_{[H]_{Q_H}} \int_{[G']_{Q_{G'}}}
    &\Gamma_{Q_{n+1}}'(H_{Q_{n+1}}(g_1')-T_{Q_{n+1}},H_{Q_{n+1}}(g_1')
    -H_{Q_{n+1}}(g_1' g_{0,1}'))\\
    & K_{(h_0, g_0') \cdot f_+, \chi}^{Q, T}(h, g')
    \eta_{n+1}(g') \rd h \rd g'.
    \end{aligned}
    \end{equation}
We calculate this as in the proof of
Theorem~\ref{thm:coarse_spectral_expansion_GL}. By the Iwasawa decomposition
it equals
    \begin{equation} \label{eq:invariant_Iwasawa}
            \begin{aligned}
    \int_{[M_{Q_H}] \times [M_{Q_{G'}}]} \int_{K_H \times K_{G'}}
    &e^{\langle -2\rho_{Q_H}, H_{Q_H}(m_H) \rangle}
    e^{\langle -2\rho_{Q_{G'}}, H_{Q_{G'}}(m') \rangle} \\
    &\Gamma'_{Q_{n+1}}(H_{Q_{n+1}}(m_1')-T_{Q_{n+1}},
    -H_{Q_{n+1}}(k_1' g_{0,1}'))\\
    & K_{(h_0, g_0') \cdot f_+,\chi}^{Q,T}(m_H k_H, m'k')
    \eta_{n+1}(m'k') \rd k' \rd k_H \rd m' \rd m_H .
    \end{aligned}
    \end{equation}


For $(m, l) \in M_{Q, +}(\bA)$, we put
    \begin{align*}
    f_{+, Q, h_0, g_0'}(m, l) =
    e^{\langle \rho_{Q_G},H_{Q_G}(m) \rangle}
    &\int_{K_H \times K'} \int_{{N_{Q_G}}(\bA)} \int_{N_{Q_{L^\vee}(\bA)}}
    ((h_0, g_0') \cdot f_+)(k_H^{-1}m n k',(l+ u)k_H)\\
    & p_Q(-H_{Q_{n+1}}(k_1' g_{0,1}')) \eta_{n+1}(k') \mu^{-1}(k_H)
    \rd n \rd  u \rd k_H \rd k'.
    \end{align*}

One check directly that for an $(m_H,m') \in [M_{Q_H}] \times [M_{Q_{G'}}]$, the expression
    \begin{align*}
       & \int_{K_H \times K_{G'}} \int_{A_{M_Q}^\infty}
    e^{\langle -2\rho_{Q_H}, H_{Q_H}(am_H) \rangle}
    e^{\langle -2\rho_{Q_{G'}}, H_{Q_{G'}}(am') \rangle}
    \Gamma'_{Q_{n+1}}(H_{Q_{n+1}}(am_1')-T_{Q_{n+1}},
    -H_{Q_{n+1}}(k_1' g_{0,1}'))\\
    & K_{(h_0, g_0') \cdot f_+,P}(am_H k_H, am'k')
    \eta_{n+1}(am'k')  \rd a \rd k_H \rd k'
    \end{align*}
    equals to
    \[
    e^{\langle \underline{\rho}_Q, H_{Q_{n+1}}(m_H) - H_{Q_{n+1}}(m_1') \rangle} c_Q e^{ \langle \underline{\rho}_Q, T \rangle} K_{f_{+,Q,h_0,g_0'},P \cap M_Q}^{M_Q}(m_H,m').
    \]

Then for any $\chi \in \fX(G)$,
 \begin{align*}
       & \int_{K_H \times K_{G'}} \int_{A_{M_Q}^\infty}
    e^{\langle -2\rho_{Q_H}, H_{Q_H}(am_H) \rangle}
    e^{\langle -2\rho_{Q_{G'}}, H_{Q_{G'}}(am') \rangle}
    \Gamma'_{Q_{n+1}}(H_{Q_{n+1}}(am_1')-T_{Q_{n+1}},
    -H_{Q_{n+1}}(k_1' g_{0,1}'))\\
    & K_{(h_0, g_0') \cdot f_+,P,\chi}(am_H k_H, am'k')
    \eta_{n+1}(am'k')  \rd a \rd k_H \rd k'
    \end{align*}
    equals to
    \[
    e^{\langle \underline{\rho}_Q, H_{Q_{n+1}}(m_H) - H_{Q_{n+1}}(m_1') \rangle} c_Q e^{ \langle \underline{\rho}_Q, T \rangle} K_{f_{+,Q,h_0,g_0'},P \cap M_Q,\chi}^{M_Q}(m_H,m').
    \]

Therefore ~\eqref{eq:invariant_Iwasawa} equals
    \[
    c_Q e^{\langle \underline{\rho}_Q,T \rangle}
    I_{\chi}^{M_Q,T}(\widetilde{f_{+, Q, h_0, g_0'}}),
    \]
where
    \[
       \widetilde{f_{+,Q,h_0,g_0'}}(m,l)= e^{ \langle -\underline{\rho}_Q,
    H_{Q_{n+1}}(m)\rangle} f_{+,Q,h_0,g_0'}(m,l), \quad (m, l) \in M_{Q, +}(\bA)
    \]

In conclusion, we have
    \[
    \eta_{n+1}(g_0') I_\chi^{T}(f_+) = \sum_{Q \in \cF}
    c_Q  e^{\langle \underline{\rho}_Q, T \rangle}
    I_{\chi}^{M_Q,T}(\widetilde{f_{+, Q, h_0, g_0'}}).
    \]
Each $I_{\chi}^{M_Q,T}(f_{+, Q, h_0, g_0'})$ is a exponential polynomial,
whose exponents are in the set $\{\underline{\rho}_P - \underline{\rho}_Q
\mid P \subset Q \}$. Since $\underline{\rho}_P$ is not trivial unless
$P = J_n$, the only term on the right hand side
that has a (possibly) nonzero purely polynomial part correspond to $Q =
J_n$. In this case $f_{+, Q, h_0, g_0'} = (h_0, g_0') \cdot f_+$ and
the pure polynomial part of the right hand side equals
$I_\chi((h_0, g_0') \cdot f_+)$.
\end{proof}

\end{document}
